% !TEX TS-program = pdflatex
\pdfoutput=1
\documentclass[11pt,a4paper]{article}

% ============================================================
% Packages
% ============================================================
\usepackage{amsmath,amsfonts,amssymb,amsthm,fullpage,booktabs}
\usepackage{float} % enables the [H] float placement specifier
\usepackage{algorithm,algpseudocode}
\usepackage[utf8]{inputenc}
\usepackage[numbers,sort&compress]{natbib}
\usepackage{hyperref}
\usepackage{xcolor}
\usepackage{enumitem}
\usepackage{microtype}
\usepackage{graphicx}
% Safe figure include: compile even if figure files are missing
\makeatletter
\newcommand{\safeincludegraphics}[2][]{%
	\begingroup
	% Make \IfFileExists search the same paths as \includegraphics
	\let\input@path\Ginput@path
	\def\@foundfile{}%
	% Try without extension first (rare), then common extensions
	\IfFileExists{#2}{\def\@foundfile{#2}}{%
		\IfFileExists{#2.pdf}{\def\@foundfile{#2.pdf}}{%
			\IfFileExists{#2.png}{\def\@foundfile{#2.png}}{%
				\IfFileExists{#2.jpg}{\def\@foundfile{#2.jpg}}{%
					\IfFileExists{#2.jpeg}{\def\@foundfile{#2.jpeg}}{%
						\IfFileExists{#2.eps}{\def\@foundfile{#2.eps}}{}%
					}%
				}%
			}%
		}%
	}%
	\ifx\@foundfile\@empty
	\fbox{\parbox[c][0.18\textheight][c]{0.9\linewidth}{\centering
			Missing figure: \texttt{\detokenize{#2}}}}%
	\else
	\includegraphics[#1]{\@foundfile}%
	\fi
	\endgroup
}
\makeatother

\usepackage{multirow}
\usepackage{subcaption}
\graphicspath{
	{./results/fig_multiobj/}
	{./results/fig_single/}
	{./results_scalability/fig_scalability/}
	{./results_robust/fig_robust/}
	{./results_param/param_study/}
}

\emergencystretch=2em
\allowdisplaybreaks

\hypersetup{
	colorlinks=true,
	linkcolor=blue,
	filecolor=magenta,
	urlcolor=cyan,
	citecolor=blue,
}

% Avoid hyperref PDF string warnings for math in bookmarks
\pdfstringdefDisableCommands{%
	\def\cO{O}%
	\def\one{1}%
	\def\J{J}%
	\def\Pmat{P}%
	\def\norm#1{||#1||}%
	\def\ip#1#2{<#1,#2>}%
}

% ============================================================
% Theorem-like
% ============================================================
\newtheorem{theorem}{Theorem}[section]
\newtheorem{lemma}[theorem]{Lemma}
\newtheorem{proposition}[theorem]{Proposition}
\newtheorem{assumption}[theorem]{Assumption}
\newtheorem{remark}[theorem]{Remark}

% ============================================================
% Macros
% ============================================================
\newcommand{\R}{\mathbb{R}}
\newcommand{\one}{\mathbf{1}}
\newcommand{\ip}[2]{\left\langle #1,#2\right\rangle}
\newcommand{\norm}[1]{\left\|#1\right\|}
\newcommand{\normtwo}[1]{\left\|#1\right\|_2}
\newcommand{\normF}[1]{\left\|#1\right\|_F}
\newcommand{\cO}{\mathcal{O}}
\newcommand{\J}{\left(\frac{1}{N}\one\one^\top\right)}
\newcommand{\Pmat}{I-\J}
\newcommand{\ProjPSD}{\Pi_{\succeq 0}}

\newcommand{\barx}{\bar x}
\newcommand{\barg}{\bar g}
\newcommand{\bars}{\bar s}

% ============================================================
% Title
% ============================================================
\title{\textbf{Distributed Gradient--Regularized Newton:\\
		Adaptive Consensus and \texorpdfstring{$\mathcal{O}(1/k^2)$}{O(1/k^2)} Global Convergence}}
\author{
	\textbf{Anonymous Authors} \\
	\textit{Institution} \\
	\texttt{email@domain.edu}
}
\date{\today}

%%%%%%%%%%%%%%%%%%%%%%%%%%%%%%%%%%%%%%%%%%%%%%%%%%%%%%%%%%%%
\begin{document}
	\maketitle
	
\begin{abstract}
	We study decentralized optimization of a convex finite-sum objective over a connected network
	of $N$ agents that communicate only with immediate neighbors.
	Second-order methods are appealing in communication-limited environments because they reduce
	the number of outer iterations, yet they are challenging to deploy in fully decentralized
	settings due to (i) the cost of maintaining curvature information across the network,
	(ii) sensitivity to stepsize selection, and (iii) network-induced disagreement and tracking errors.
	Motivated by stepsize-free regularized Newton schemes \citep{mishchenko2023rn},
	we propose \textsc{DisGrem}, a decentralized gradient--regularized Newton method that combines
	gradient tracking, Hessian tracking with PSD projection, and a two-stage mixing
	(pre-mixing before the local Newton solve, post-mixing after the update).
	A user-chosen scaling factor $M>0$ amplifies the vanishing regularization
	$\lambda_{i,k}=\sqrt{M\|\tilde g_{i,k}\|}$, enlarging the admissible heterogeneity regime
	while preserving $\lambda_{i,k}\to 0$ as the gradient vanishes.
	
	Our analysis introduces a \emph{reference-step} construction that resolves the
	non-commutativity between network averaging and matrix inversion, yielding an
	\emph{inexact regularized Newton} characterization of the average iterate.
	To quantify persistent curvature mismatch across agents, we introduce a bounded Hessian-heterogeneity
	constant over the relevant level set; it enters our network-dependent constants and is mild in near-i.i.d.\ regimes.
	Under convexity, Lipschitz Hessians, bounded level sets, and sufficiently accurate mixing,
	we establish $\cO(1/k^2)$ functional decay and $\cO(\varepsilon^{-1})$ iteration complexity
	to reach $\|\nabla f(\bar x_k)\|\le \varepsilon$ \emph{above a heterogeneity-dependent accuracy floor}
	(and recover the classical centralized guarantees in homogeneous/zero-bias regimes).
	Moreover, under a logarithmic accuracy schedule we obtain
	$\cO(\varepsilon^{-1}\log(\varepsilon^{-1}))$ total neighbor-communication rounds.
	
	We also develop practical variants: \textsc{CeDisGrem} compresses Hessian corrections
	to reduce communication by up to $50\%$, and \textsc{AdaDisGrem} adapts $M$ online.
	Extensive numerical experiments on 15 benchmark objectives, including ill-conditioned
	quadratics, log-sum-exp, real-data logistic regression, and non-convex landscapes,
	show that the \textsc{DisGrem} family converges to near machine precision $5$--$8\times$ faster
	than the best first-order baseline (EXTRA) and robustly handles weakly connected networks
	where gradient-tracking first-order methods fail entirely.
\end{abstract}
	
	%%%%%%%%%%%%%%%%%%%%%%%%%%%%%%%%%%%%%%%%%%%%%%%%%%%%%%%%%%%%
	\section{Introduction}
	%%%%%%%%%%%%%%%%%%%%%%%%%%%%%%%%%%%%%%%%%%%%%%%%%%%%%%%%%%%%
	Decentralized optimization addresses problems of the form
	\begin{equation}\label{eq:intro_prob}
		\min_{x\in\R^d} f(x)\;:=\;\frac{1}{N}\sum_{i=1}^N f_i(x),
	\end{equation}
	where each $f_i$ is stored privately at node $i$ of a connected network and nodes communicate only with neighbors.
	A principal bottleneck is \emph{communication}: algorithms must trade off local computation
	against costly neighbor exchanges.
	First-order decentralized methods may require many iterations on ill-conditioned problems, amplifying communication costs.
	Second-order information can reduce outer iterations, but decentralized second-order methods face major obstacles:
	curvature information is expensive to maintain and global convergence often requires careful stepsize control or line-search.
	
	A stepsize-free perspective is provided by the regularized Newton method of \citet{mishchenko2023rn},
	which uses a vanishing regularization scale $\lambda_k\asymp \sqrt{\|g_k\|}$ and achieves the optimal functional rate $\cO(1/k^2)$
	under convexity and Lipschitz Hessians without manual stepsize tuning.
	This motivates: \emph{can we obtain the same outer-iteration complexity in a decentralized network while controlling disagreement/tracking errors?}
	
\noindent\textbf{Contributions.}
We propose and analyze \textsc{DisGrem} (Distributed Gradient-Regularized Newton Method),
which combines gradient tracking, Hessian tracking with PSD projection, and two-stage
mixing (pre-mixing before the local Newton solve, post-mixing after the update).
Our specific contributions are:
\begin{itemize}[leftmargin=1.2em,itemsep=2pt]
	\item \textbf{Algorithm design.}
	We formulate a stepsize-free decentralized second-order method based on vanishing regularization
	$\lambda_{i,k}=\sqrt{M\|\tilde g_{i,k}\|}$.
	Two-stage mixing decouples the network averaging accuracy from the Newton solve,
	and PSD projection of tracked Hessians maintains positive semidefiniteness throughout.
	
	\item \textbf{Reference-step framework.}
	We introduce a reference-step construction to control the non-commutativity between network averaging
	and matrix inversion, yielding an \textbf{inexact regularized Newton} characterization for the average iterate.
	
	\item \textbf{Hessian heterogeneity analysis.}
	We identify the obstruction from persistent inter-agent Hessian mismatch and quantify it through
	a bounded heterogeneity constant $\sigma_H$ on the relevant level set; see~\eqref{eq:sigma_H_def}.
	
	\item \textbf{Convergence guarantees.}
	Under stated assumptions we prove $\Phi_k=\cO(1/k^2)$ functional decay and
	\textbf{$\cO(\varepsilon^{-1})$} iteration complexity to reach $\|\nabla f(\bar x_k)\|\le \varepsilon$
	for any $\varepsilon$ above the explicit floor in~\eqref{eq:eps_floor},
	and \textbf{$\cO(\varepsilon^{-1}\log(\varepsilon^{-1}))$} total mixing rounds under a logarithmic accuracy schedule.
	
	\item \textbf{Practical variants.}
	We develop \textbf{\textsc{CeDisGrem}}, which applies top-$r$ low-rank compression to Hessian correction messages
	and reduces per-iteration communication by 30--50\% with negligible iteration overhead;
	and \textbf{\textsc{AdaDisGrem}}, which adapts the scaling factor $M$ online to track local curvature changes.
	
	\item \textbf{Numerical evidence.}
	Experiments on 15 objectives across well-conditioned and ill-conditioned convex problems, real-data logistic regression,
	and non-convex landscapes show that the \textsc{DisGrem} family consistently converges to near machine precision while
	first-order baselines stagnate on weakly connected networks.
	It achieves $5$--$8\times$ fewer iterations than EXTRA and 100\% success rate in Monte-Carlo robustness trials
	where gradient-tracking baselines fail.
\end{itemize}
	
	%%%%%%%%%%%%%%%%%%%%%%%%%%%%%%%%%%%%%%%%%%%%%%%%%%%%%%%%%%%%
	\section{Related Work}
	%%%%%%%%%%%%%%%%%%%%%%%%%%%%%%%%%%%%%%%%%%%%%%%%%%%%%%%%%%%%
	
	\noindent\textbf{First-order decentralized methods.}
	Decentralized (sub)gradient descent dates to \citet{nedic2009subgradient}.
	The bias of constant-stepsize DGD is corrected by exact methods:
	EXTRA~\citep{shi2015extra} uses a two-step update with a carefully chosen pair of matrices,
	and exact diffusion~\citep{yuan2019exactdiffusion} reformulates the update in a primal-dual
	framework.
	Gradient tracking~\citep{nedic2017diging} achieves exact convergence
	under a single doubly stochastic matrix by maintaining a running sum of gradient increments.
	These methods achieve $\cO(1/k)$ or linear convergence for strongly convex objectives,
	but their iteration complexity depends on $(1-\rho)^{-1}$ or $(1-\rho)^{-2}$, which becomes
	prohibitive on weakly connected networks.
	Fast distributed gradient methods~\citep{jakovetic2014fast} improve over pure gradient descent
	but remain first-order.
	
	\noindent\textbf{Second-order and quasi-Newton decentralized methods.}
	Network Newton~\citep{mokhtari2015networknewton_arxiv} approximates the global Newton direction
	by a truncated power series of the Hessian, using multi-hop communication.
	Newton tracking~\citep{mokhtari2020newtontracking} integrates Newton-type directions into the
	tracking framework.
	Decentralized quasi-Newton methods, including DQM~\citep{eisen2017dqn}, replace the exact
	Hessian with L-BFGS-type approximations.
	All of these methods require explicit stepsize tuning and do not provide global convergence
	guarantees matching the centralized Newton rate.
	The challenge of Hessian heterogeneity has not been explicitly characterized in prior work.
	
	\noindent\textbf{Regularized Newton and cubic regularization.}
	The stepsize-free perspective adopted here is inspired by the centralized regularized Newton
	method of \citet{mishchenko2023rn}, which uses $\lambda_k\asymp\sqrt{\|\nabla f(x_k)\|}$
	and achieves the tight $\cO(1/k^2)$ functional rate under convexity and Lipschitz Hessians.
	This builds on cubic-regularization ideas of \citet{nesterov2006cubic,cartis2011arc},
	which also achieve $\cO(1/k^2)$ via an adaptive cubic model.
	Our paper is the first, to our knowledge, to show that the centralized $\cO(1/k^2)$
	rate can be preserved in a decentralized network via Hessian/gradient tracking and
	two-stage mixing, and the first to provide a transparent treatment of the Hessian
	heterogeneity obstruction.
	
	%%%%%%%%%%%%%%%%%%%%%%%%%%%%%%%%%%%%%%%%%%%%%%%%%%%%%%%%%%%%
	\section{Preliminaries and Assumptions}\label{sec:prelim}
	%%%%%%%%%%%%%%%%%%%%%%%%%%%%%%%%%%%%%%%%%%%%%%%%%%%%%%%%%%%%
	\subsection{Network model and notation}
	Communication uses a symmetric, nonnegative, doubly stochastic mixing matrix $W\in\R^{N\times N}$:
	$W_{ij}\ge0$, $W\one=\one$, and $\one^\top W=\one^\top$.
	Define the averaging projector $\J=\frac{1}{N}\one\one^\top$ and $\Pmat=I-\J$.
	Let $\rho := \normtwo{W-\J}\in[0,1)$, strictly less than $1$ for a connected undirected network.
	For stacked vectors $Z=[z_1;\dots;z_N]\in\R^{Nd}$ we use the Euclidean norm $\|Z\|$.
	We write $\bar z:=\frac1N\sum_{i=1}^N z_i$ for averages.
	
	\subsection{Standing assumptions}
\begin{assumption}\label{ass:standing}
	\begin{enumerate}\itemsep2pt
		\item \textbf{Convexity and smoothness.}
		Each $f_i$ is convex and twice continuously differentiable.
		Its gradient is $L_1$-Lipschitz and its Hessian is $L_2$-Lipschitz:
		\[
		\norm{\nabla f_i(x)-\nabla f_i(y)}\le L_1\norm{x-y},
		\qquad
		\normtwo{\nabla^2 f_i(x)-\nabla^2 f_i(y)}\le L_2\norm{x-y}.
		\]
		\item \textbf{PSD local curvature.}
		$\nabla^2 f_i(x)\succeq0$ for all $x$ and all $i$.
		\item \textbf{Connected network.}
		$W$ is symmetric, doubly stochastic, and $\rho=\normtwo{W-\J}<1$.
		\item \textbf{Bounded level set (iterates).}
		There exist $D>0$ and $x_\star\in\arg\min f$ such that for all $k\ge 0$ and all $i$,
		\begin{equation}\label{eq:bounded_levelset}
			\norm{x_{i,k}-x_\star}\le D.
		\end{equation}
		Consequently, $\norm{\bar x_k-x_\star}\le D$ and $\norm{\tilde x_{i,k}-x_\star}\le D$
		for all $i,k$ because $\bar x_k$ and $\tilde x_{i,k}$ are convex combinations of $\{x_{j,k}\}$.
		\item \textbf{Bounded Hessians on the relevant set.}
		There exists $M_H\ge0$ such that $\normtwo{\nabla^2 f(x)}\le M_H$ for all $x$ with $\norm{x-x_\star}\le D$.
		\item \textbf{Regularization scaling.} The user-chosen scaling factor satisfies $M\ge L_2$.
	\end{enumerate}
\end{assumption}
	
	\begin{remark}[Interpretation of Assumption~\ref{ass:standing}(6)]
		Assumption~\ref{ass:standing}(6) is a \emph{boundedness} condition on curvature mismatch over the relevant level set.
		It is automatically satisfied whenever $\nabla^2 f_i$ are bounded on $\{x:\ \norm{x-x_\star}\le D\}$, and it becomes small in near-i.i.d.\ regimes.
		Unlike a vanishing condition, it does \emph{not} force local Hessians to coincide at stationary points; rather, it quantifies the irreducible mismatch that must be controlled by mixing and regularization.
	\end{remark}
	
	% ============================================================
	% Derived heterogeneity constant (finite on the bounded level set)
	% ============================================================
	Define the (possibly zero) gradient-heterogeneity constant
	\begin{equation}\label{eq:sigma_g_def}
		\sigma_g := \sup\Big\{\, \norm{((\Pmat)\otimes I_d)\,\nabla F(\one\otimes x)}\;:\; \norm{x-x_\star}\le D\,\Big\}\in[0,\infty),
	\end{equation}
	which is finite under Assumption~\ref{ass:standing}(1) and (4) because $\nabla F(\one\otimes x)$ is continuous on a compact set.
	In interpolation/shared-minimizer regimes one typically has $\sigma_g=0$, while for general heterogeneous data this constant may be positive and induces an unavoidable accuracy floor.
	
	
	%%%%%%%%%%%%%%%%%%%%%%%%%%%%%%%%%%%%%%%%%%%%%%%%%%%%%%%%%%%%
	\section{Algorithm}\label{sec:alg}
	%%%%%%%%%%%%%%%%%%%%%%%%%%%%%%%%%%%%%%%%%%%%%%%%%%%%%%%%%%%%
	
	\subsection{Design motivation}\label{ssec:design}
	
	Three interacting challenges must be resolved to lift regularized Newton to the decentralized setting.
	
	\textbf{(i) Averaging vs.\ inversion.}
	In the centralized regularized Newton method of \citet{mishchenko2023rn} the Newton system involves
	$(\nabla^2 f(\bar x_k)+\lambda_k I)^{-1}$, which is an invertible function of the \emph{average} gradient.
	In a network each agent solves its own local system, so the true average step
	$\bar s_k = \frac1N\sum_i s_{i,k} \ne (\frac1N\sum_i A_{i,k})^{-1}(\frac1N\sum_i g_{i,k})$
	in general: averaging and inversion do not commute.
	Pre-mixing reduces disagreement in the inputs $(x_{i,k}, g_{i,k}, H_{i,k})$ before the solve,
	making individual systems closer to the average system and bounding the resulting
	\emph{step dispersion} (Section~\ref{sec:theory}, Lemma~\ref{lem:step-disp}).
	
	\textbf{(ii) Hessian maintenance and positivity.}
	Tracker dynamics propagate Hessian increments $\nabla^2 f_j(x_{j,k+1})-\nabla^2 f_j( x_{j,k})$
	through the network.
	Floating-point errors and network delays may produce non-symmetric or indefinite tracked Hessians,
	which would make the regularized Newton system ill-conditioned.
	We therefore symmetrize and project $\widehat H_{i,k+1}$ onto the PSD cone in Frobenius norm
	before storing $H_{i,k+1}$.
	This projection is non-expansive (Remark~\ref{rem:psd-nonexp}) and introduces an error
	$\hat e_{i,k}:=\normF{\tilde H_{i,k}-\nabla^2 f_i(\tilde x_{i,k})}$ that is absorbed into
	the regularization via the augmented shift $\lambda_{i,k}+\hat e_{i,k}$,
	guaranteeing a well-posed local solve for any tracked Hessian.
	
	\textbf{(iii) Stepsize-free regularization.}
	Setting $\lambda_{i,k}=\sqrt{M\|\tilde g_{i,k}\|}$ mirrors the centralized choice
	$\lambda_k\asymp\sqrt{\|\nabla f(\bar x_k)\|}$, which achieves $\cO(1/k^2)$ without a line search
	or manual stepsize.
	The user constant $M\ge L_2$ scales the regularizer uniformly: larger $M$ produces larger shifts,
	which damp Newton steps, tolerate wider heterogeneity, and improve stability on
	ill-conditioned or non-convex problems (Section~\ref{ssec:param});
	once the iterates enter a region where $\|\nabla f\|\to 0$, the regularizer vanishes automatically.
	
	\subsection{\textsc{DisGrem}: full algorithm}\label{ssec:disgrem}
	
	Each node $i$ stores a primal variable $x_{i,k}\in\R^d$, a gradient tracker $g_{i,k}\in\R^d$,
	and a Hessian tracker $H_{i,k}\in\R^{d\times d}$.
	$\ProjPSD(\cdot)$ denotes the Euclidean (Frobenius) projection onto the PSD cone.
	
\begin{algorithm}[H]
	\caption{Decentralized Gradient--Regularized Newton (\textsc{DisGrem}, v3)}
	\label{alg:disgrem}
	\begin{algorithmic}[1]
		\State \textbf{Input:} $\{x_{i,0}\}_{i=1}^N$, mixing matrix $W$, Lipschitz constants $L_1,L_2>0$, scaling factor $M\ge L_2$.
		\State \textbf{Initialize:}
		$g_{i,0}\gets \nabla f_i(x_{i,0})$,
		$H_{i,0}\gets \nabla^2 f_i(x_{i,0})$ (stored as a symmetric matrix).
		\For{$k=0,1,2,\dots$}
		\State \textbf{(A) Pre-mixing:} choose $\tau_k\ge 1$ and set
		\[
		\tilde x_{i,k}\gets \sum_{j=1}^N [W^{\tau_k}]_{ij}\,x_{j,k},\quad
		\tilde g_{i,k}\gets \sum_{j=1}^N [W^{\tau_k}]_{ij}\,g_{j,k},\quad
		\tilde H_{i,k}\gets \sum_{j=1}^N [W^{\tau_k}]_{ij}\,H_{j,k}.
		\]
		\State \textbf{(B) Local computation (parallel for each $i$):}
		\[
		\lambda_{i,k} \gets \sqrt{M\,\norm{\tilde g_{i,k}}}.
		\]
		\Statex\hspace{1.6em}Let $\mathrm{sym}(A):=\tfrac12(A+A^\top)$ and define the PSD-stabilizer
		\[
		\delta_{i,k}\ :=\ \max\Big\{0,\ -\lambda_{\min}\big(\mathrm{sym}(\tilde H_{i,k})\big)\Big\}\ \ \ge 0.
		\]
		\Statex\hspace{1.6em}\textbf{Solve}
		\[
		\big(\mathrm{sym}(\tilde H_{i,k})+(\lambda_{i,k}+\delta_{i,k})I\big)\,s_{i,k}=-\tilde g_{i,k}
		\]
		\Statex\hspace{1.6em}(set $s_{i,k}=0$ if $\tilde g_{i,k}=0$); set $y_{i,k+1}\gets \tilde x_{i,k}+s_{i,k}$.
		\State \textbf{(C) Post-mixing:} choose $t_k\ge 1$ and set
		\[
		x_{i,k+1}\gets \sum_{j=1}^N [W^{t_k}]_{ij}\,y_{j,k+1}.
		\]
		\State \textbf{(D) Tracker updates (increment form):}
		\[
		g_{i,k+1}\gets \sum_{j=1}^N [W^{t_k}]_{ij}\Big(\tilde g_{j,k}+\nabla f_j(x_{j,k+1})-\nabla f_j(x_{j,k})\Big),
		\]
		\[
		\widehat H_{i,k+1}\gets \sum_{j=1}^N [W^{t_k}]_{ij}\Big(\tilde H_{j,k}+\nabla^2 f_j(x_{j,k+1})-\nabla^2 f_j(x_{j,k})\Big),
		\qquad
		H_{i,k+1}\gets \mathrm{sym}(\widehat H_{i,k+1}).
		\]
		\EndFor
	\end{algorithmic}
\end{algorithm}
	
\subsection{Communication cost per iteration}\label{ssec:comm}

Each outer iteration of \textsc{DisGrem} involves $(\tau_k + t_k)$ rounds of neighbor communication.
In each round every agent $i$ sends messages to and receives messages from its $|\mathcal N_i|$ neighbors.
In \textsc{DisGrem}, the natural payload exchanged per neighbor per round consists of:
\begin{itemize}[leftmargin=1.2em,itemsep=1pt]
	\item two $d$-dimensional vectors (e.g., mixing of $(x,g)$ in pre-mixing, or mixing of $(y,g)$ in post-mixing): $2d$ floats,
	\item one $d\times d$ symmetric matrix (Hessian tracker or Hessian increment): $\tfrac{d(d+1)}{2}$ floats.
\end{itemize}
Equivalently, the two vectors can be concatenated into a single $2d$-vector; this does not change the byte complexity.

Let $n_v$ denote the number of $d$-vectors transmitted per round (for \textsc{DisGrem}, $n_v=2$).
With double precision (8 bytes per float), the \emph{sent} message volume per agent per outer iteration is
\begin{equation}\label{eq:comm_cost_iter}
	C_{\mathrm{iter}}^{\mathrm{send}}(k)
	= (\tau_k+t_k)\cdot |\mathcal N_i|\cdot \Bigl(n_v d + \tfrac{d(d+1)}{2}\Bigr)\cdot 8 \ \ \text{bytes}.
\end{equation}
If one counts both sending and receiving volume, then
$C_{\mathrm{iter}}^{\mathrm{send+recv}}(k)=2\,C_{\mathrm{iter}}^{\mathrm{send}}(k)$.

For the default setting $N=10$, $d=10$, and the random geometric graph with $r_0=0.5$ giving average degree
$|\mathcal N_i|\approx 3.5$, taking $\tau_k=t_k=3$ and $n_v=2$ yields
\[
C_{\mathrm{iter}}^{\mathrm{send}}\approx 6\times 3.5\times\Bigl(20+55\Bigr)\times 8
\approx 12.6\,\text{kB}
\quad \text{per agent per outer iteration}.
\]
This back-of-the-envelope estimate is consistent with the empirical \texttt{commCost} magnitudes reported in
Table~\ref{tab:main} (up to graph-degree variability and implementation details such as message aggregation and bookkeeping).
	
	\subsection{Practical variants}\label{ssec:variants}
	
\noindent\textbf{\textsc{CeDisGrem}: compressed Hessian communication.}
The dominant communication cost of \textsc{DisGrem} is the $\frac{d(d+1)}{2}$-dimensional
Hessian increment $\nabla^2 f_j(x_{j,k+1})-\nabla^2 f_j(x_{j,k})$.
\textsc{CeDisGrem} replaces each symmetric Hessian increment
$\Delta H_{j,k}:=\nabla^2 f_j(x_{j,k+1})-\nabla^2 f_j(x_{j,k})$
by a rank-$r$ symmetric truncation based on the eigen-decomposition.
Let $\Delta H_{j,k}=V\operatorname{diag}(\mu_1,\dots,\mu_d)V^\top$ with
$|\mu_1|\ge \cdots \ge |\mu_d|$ and $V=[v_1,\dots,v_d]$ orthogonal.
Define the truncated approximation
\[
\hat\Delta H_{j,k}:=\sum_{\ell=1}^r \mu_\ell v_\ell v_\ell^\top
= V_r\Lambda_r V_r^\top,
\]
where $V_r=[v_1,\dots,v_r]\in\R^{d\times r}$ and $\Lambda_r=\operatorname{diag}(\mu_1,\dots,\mu_r)\in\R^{r\times r}$.
This $\hat\Delta H_{j,k}$ is the best rank-$r$ symmetric approximation of $\Delta H_{j,k}$ in Frobenius norm
(equivalently, the best rank-$r$ approximation for a symmetric matrix is obtained by keeping the $r$
eigenpairs with largest absolute eigenvalues).
Thus, transmitting $(V_r,\Lambda_r)$ costs $r(d+1)$ floats instead of $\frac{d(d+1)}{2}$ floats per neighbor per round.

The truncation error satisfies
\[
\normF{\Delta H_{j,k}-\hat\Delta H_{j,k}}
=\left(\sum_{\ell=r+1}^d \mu_\ell^2\right)^{1/2}
\le |\mu_{r+1}(\Delta H_{j,k})|\sqrt{d-r},
\]
where $\mu_{r+1}(\Delta H_{j,k})$ denotes the $(r+1)$-th largest absolute eigenvalue.
This error is absorbed into $\hat e_{i,k}$ and thus into the regularization automatically.
Empirically, $r=\lceil d/5\rceil$ reduces communication by 30--50\% with negligible iteration overhead
(Table~\ref{tab:main}).
	
	\noindent\textbf{\textsc{AdaDisGrem}: adaptive scaling factor.}
	The default choice $M=\mathrm{const}$ requires prior knowledge of $L_2$.
	\textsc{AdaDisGrem} maintains a running estimate $\hat M_k$ of the local Hessian Lipschitz
	constant using the secant condition:
	\[
	\hat M_k \gets \max\!\Bigl(\hat M_{k-1},\;
	\frac{\normtwo{\nabla^2 f_i(x_{i,k})-\nabla^2 f_i(x_{i,k-1})}}{\norm{x_{i,k}-x_{i,k-1}}}\Bigr),
	\]
	and sets $\lambda_{i,k}=\sqrt{\hat M_k\|\tilde g_{i,k}\|}$.
	This heuristic removes the need to specify $L_2$ in practice; a full theoretical treatment of the adaptive rule is left for future work.
	
	\subsection{Practical guidelines for parameter selection}\label{ssec:practical}
	
	\begin{itemize}[leftmargin=1.2em,itemsep=2pt]
		\item \textbf{Scaling factor $M$.}
		Theorem~\ref{thm:main} requires $M\ge L_2$.
		In practice, setting $M = c \cdot \hat L$, where $\hat L$ is a rough estimate of the
		gradient Lipschitz constant and $c\in[0.5, 2.0]$, works well across all tested problems.
		Our parameter study (Section~\ref{ssec:param}) shows convergence is largely insensitive
		to the exact value of $M$ within this range.
		When $L_2$ is unknown, use \textsc{AdaDisGrem}.
		\item \textbf{Consensus rounds $N_c$.}
		The theory requires $\tau_k, t_k = \cO(\log 1/\|\tilde{\bar g}_k\|)$, which grows
		logarithmically.  In practice, a fixed $N_c = 3$--$5$ rounds per stage is sufficient
		for most networks with $\rho \le 0.95$.
		For strongly sparse networks ($\rho > 0.97$), increasing $N_c$ to $8$--$10$
		significantly improves stability.
		\item \textbf{Hessian compression rank $r$ (\textsc{CeDisGrem}).}
		Set $r = \lceil 0.2d \rceil$ as a default; increase $r$ if convergence slows.
		For diagonal-dominant Hessians (e.g.\ $\ell_2$-regularized losses), even $r=1$
		is often sufficient.
	\end{itemize}
	
%%%%%%%%%%%%%%%%%%%%%%%%%%%%%%%%%%%%%%%%%%%%%%%%%%%%%%%%%%%%
\section{Theory and Convergence Analysis}\label{sec:theory}
%%%%%%%%%%%%%%%%%%%%%%%%%%%%%%%%%%%%%%%%%%%%%%%%%%%%%%%%%%%%
Throughout, Assumption~\ref{ass:standing} holds.
All detailed proofs are provided in Appendices~\ref{app:dispersion}--\ref{app:proofs}.

\subsection{Roadmap of the analysis}\label{ssec:roadmap}
Our theory is organized into four layers.
\begin{itemize}[leftmargin=1.2em,itemsep=2pt]
	\item \textbf{Layer I (identities):} averaging is preserved by mixing; the tracker averages match the true averaged gradient/Hessian at the current local iterates (Lemma~\ref{lem:avg-pres}--\ref{lem:avg-track-h}).
	\item \textbf{Layer II (reference step):} we build a \emph{reference} regularized Newton step based on averaged pre-mixed quantities and control the non-commutativity between averaging and inversion via a resolvent bound (Lemma~\ref{lem:step-disp}).
	\item \textbf{Layer III (inexact RN condition):} combining the bridge bounds and the reference-step control yields an \emph{inexact regularized Newton} condition for the average iterate (Proposition~\ref{prop:inexact-rn}).
	\item \textbf{Layer IV (rates):} the inexact RN condition implies $\cO(1/k^2)$ functional decay and $\cO(\varepsilon^{-1})$ iteration complexity (Theorem~\ref{thm:main}), and an explicit logarithmic mixing schedule yields $\cO(\varepsilon^{-1}\log(\varepsilon^{-1}))$ total communication rounds (Theorem~\ref{thm:comm}).
\end{itemize}
Under strong convexity, we further derive a local $Q$-superlinear ($3/2$-order) rate for the gradient norm (Theorem~\ref{thm:sc-superlinear}).

% ============================================================
\subsection{Basic inequalities}
% ============================================================
\begin{lemma}[Third-order Taylor bound]\label{lem:taylor3}
	Let $h:\R^d\to\R$ have $L_2$-Lipschitz Hessian. Then for all $x,s$,
	\[
	h(x+s)\le h(x)+\ip{\nabla h(x)}{s}+\frac12 s^\top\nabla^2 h(x)s+\frac{L_2}{6}\norm{s}^3.
	\]
\end{lemma}

\begin{lemma}[Consensus contraction]\label{lem:consensus}
	For any $z\in\R^N$ and integer $t\ge 1$,
	\[
	\normtwo{(W^t-\J)z}\le \rho^t\normtwo{z}.
	\]
	Consequently, for any stacked $Z\in\R^{Nd}$,
	\[
	\norm{((W^t-\J)\otimes I_d)Z}\le \rho^t\norm{Z}.
	\]
\end{lemma}

\begin{lemma}[Smooth convex gap--gradient inequality]\label{lem:gap-grad}
	Let $h$ be convex with $L$-Lipschitz gradient. Then for any minimizer $x_\star$,
	\[
	\norm{\nabla h(x)}^2\le 2L\big(h(x)-h(x_\star)\big).
	\]
\end{lemma}

% ============================================================
\subsection{Average identities and reference-step construction}
% ============================================================
Define averages $\bar x_k:=\frac1N\sum_i x_{i,k}$ and likewise for other variables.
Define the pre-mixed averages
\[
\tilde{\bar x}_k:=\frac1N\sum_{i=1}^N\tilde x_{i,k},\quad
\tilde{\bar g}_k:=\frac1N\sum_{i=1}^N\tilde g_{i,k},\quad
\tilde{\bar H}_k:=\frac1N\sum_{i=1}^N\tilde H_{i,k}.
\]
Let $\bar s_k:=\frac1N\sum_{i=1}^N s_{i,k}$ denote the average step.

\begin{lemma}[Average preservation]\label{lem:avg-pres}
	For all $k$, $\tilde{\bar x}_k=\bar x_k$ and $\bar x_{k+1}=\bar x_k+\bar s_k$.
\end{lemma}

\begin{lemma}[Average tracking identity for gradients]\label{lem:avg-track}
	Assume Algorithm~\ref{alg:disgrem} uses the tracker update in step \textbf{(D)} and is initialized by
	$g_{i,0}=\nabla f_i(x_{i,0})$.
	Then for all $k\ge 0$,
	\[
	\bar g_k:=\frac1N\sum_{i=1}^N g_{i,k}=\frac1N\sum_{i=1}^N \nabla f_i(x_{i,k}),
	\qquad
	\tilde{\bar g}_k=\bar g_k.
	\]
\end{lemma}

\begin{lemma}[Average tracking identity for Hessians]\label{lem:avg-track-h}
	Assume Algorithm~\ref{alg:disgrem} uses the Hessian tracker update in step \textbf{(D)} and is initialized by
	$H_{i,0}=\nabla^2 f_i(x_{i,0})$.
	Then for all $k\ge 0$,
	\[
	\bar H_k:=\frac1N\sum_{i=1}^N H_{i,k}=\frac1N\sum_{i=1}^N \nabla^2 f_i(x_{i,k}),
	\qquad
	\tilde{\bar H}_k=\bar H_k.
	\]
\end{lemma}

\subsection{Well-posed local solves and a reference step}
Let $\mathrm{sym}(A):=\tfrac12(A+A^\top)$.
Algorithm~\ref{alg:disgrem} sets
\[
\lambda_{i,k}:=\sqrt{M\|\tilde g_{i,k}\|},
\qquad
\delta_{i,k}:=\max\{0,-\lambda_{\min}(\mathrm{sym}(\tilde H_{i,k}))\},
\qquad
\tilde\lambda_{i,k}:=\lambda_{i,k}+\delta_{i,k}.
\]
Define $A_{i,k}:=\mathrm{sym}(\tilde H_{i,k})+\tilde\lambda_{i,k}I\succeq \tilde\lambda_{i,k}I$ and
$s_{i,k}=-A_{i,k}^{-1}\tilde g_{i,k}$.

Define averaged quantities
\[
\tilde{\bar\lambda}_k:=\frac1N\sum_{i=1}^N\tilde\lambda_{i,k},
\qquad
A_k^{\mathrm{ref}}:=\mathrm{sym}(\tilde{\bar H}_k)+\tilde{\bar\lambda}_k I,
\qquad
s_k^{\mathrm{ref}}:=-(A_k^{\mathrm{ref}})^{-1}\tilde{\bar g}_k.
\]

\begin{lemma}[Scaled regularization dominates step size]\label{lem:M-vs-s}
	For every $i,k$, the local step satisfies
	\[
	M\norm{s_{i,k}}\le \lambda_{i,k},
	\qquad\text{and hence}\qquad
	L_2\norm{s_{i,k}}\le \lambda_{i,k}.
	\]
\end{lemma}

\begin{lemma}[Average regularization regime]\label{lem:avg-regime}
	For all $k$,
	\[
	L_2\norm{\bar s_k}\le \frac1N\sum_{i=1}^N \lambda_{i,k}\le \tilde{\bar\lambda}_k.
	\]
\end{lemma}

\subsection{Bridge bounds (no heterogeneity constants)}
Define the post-mixing dispersions
\[
D_X(k):=\sqrt{\frac1N\sum_{i=1}^N \|x_{i,k}-\bar x_k\|^2},
\qquad
D_H(k):=\sqrt{\frac1N\sum_{i=1}^N \|\nabla^2 f_i(x_{i,k})-\nabla^2 f_i(\bar x_k)\|_F^2}.
\]
(We will use $D_X(k)$ as the main driver; $D_H(k)\le \sqrt d\,L_2 D_X(k)$ by Lipschitzness.)

\begin{lemma}[Gradient bridge]\label{lem:bridge-grad}
	For every $k$,
	\[
	\norm{\nabla f(\bar x_k)-\tilde{\bar g}_k}\le L_1\,D_X(k).
	\]
\end{lemma}

\begin{lemma}[Hessian bridge]\label{lem:bridge-hess}
	For every $k$,
	\[
	\normtwo{\nabla^2 f(\bar x_k)-\mathrm{sym}(\tilde{\bar H}_k)}\le L_2\,D_X(k).
	\]
\end{lemma}

\begin{lemma}[Stabilizer is controlled by tracking/consensus errors]\label{lem:delta_control}
	For every $k$ and every $i$,
	\[
	0\le \delta_{i,k}\le \Big\|\mathrm{sym}(\tilde H_{i,k})-\nabla^2 f(\bar x_k)\Big\|_2
	\le \Big\|\mathrm{sym}(\tilde H_{i,k})-\mathrm{sym}(\tilde{\bar H}_k)\Big\|_2
	+\Big\|\mathrm{sym}(\tilde{\bar H}_k)-\nabla^2 f(\bar x_k)\Big\|_2.
	\]
	Consequently,
	\[
	\frac1N\sum_{i=1}^N\delta_{i,k}
	\ \le\ \Delta_k^H + L_2 D_X(k),
	\qquad
	\Delta_k^H:=\frac1N\sum_{i=1}^N\big\|\mathrm{sym}(\tilde H_{i,k})-\mathrm{sym}(\tilde{\bar H}_k)\big\|_2.
	\]
\end{lemma}

\subsection{Step dispersion via a resolvent identity}
Define dispersion measures
\[
\Delta^g_k:=\frac1N\sum_{i=1}^N \norm{\tilde g_{i,k}-\tilde{\bar g}_k},\quad
\Delta^H_k:=\frac1N\sum_{i=1}^N \normtwo{\mathrm{sym}(\tilde H_{i,k})-\mathrm{sym}(\tilde{\bar H}_k)},\quad
\Delta^\lambda_k:=\frac1N\sum_{i=1}^N \big|\tilde\lambda_{i,k}-\tilde{\bar\lambda}_k\big|,
\]
and $\underline\lambda_k:=\min_{1\le i\le N}\tilde\lambda_{i,k}$.

\begin{lemma}[Step dispersion via resolvent identity]\label{lem:step-disp}
	For every $k$,
	\[
	\norm{\bar s_k-s^{\mathrm{ref}}_k}
	\le \frac{1}{\underline\lambda_k}\Delta^g_k
	+\frac{\norm{\tilde{\bar g}_k}}{\underline\lambda_k\,\tilde{\bar\lambda}_k}\big(\Delta^H_k+\Delta^\lambda_k\big).
	\]
\end{lemma}

% ============================================================
\subsection{Inexact regularized Newton condition for the average step}
% ============================================================
Set $\eta:=1/12$ and define
\[
g_k:=\nabla f(\bar x_k),
\qquad
\lambda_k:=\tilde{\bar\lambda}_k,
\qquad
r_k := (\nabla^2 f(\bar x_k)+\lambda_k I)\bar s_k+g_k.
\]

\begin{proposition}[Inexact RN condition for the average step]\label{prop:inexact-rn}
	Assume the following hold at iteration $k$ for some $\eta\in(0,1)$:
	\begin{enumerate}[leftmargin=1.2em,itemsep=2pt]
		\item \textbf{Dispersion control:}
		\[
		\norm{\bar s_k-s^{\mathrm{ref}}_k}\le \frac{\eta}{8}\cdot \frac{\lambda_k}{M_H+\lambda_k}\,\norm{\bar s_k}.
		\]
		\item \textbf{Gradient bridge accuracy:}
		\[
		L_1 D_X(k)\le \frac{\eta}{8}\,\lambda_k\norm{\bar s_k}.
		\]
		\item \textbf{Hessian bridge accuracy:}
		\[
		L_2 D_X(k)\le \frac{\eta}{8}\,\lambda_k.
		\]
	\end{enumerate}
	Then
	\[
	\norm{r_k}\le \eta\,\lambda_k\norm{\bar s_k},
	\qquad\text{and}\qquad
	L_2\norm{\bar s_k}\le \lambda_k.
	\]
\end{proposition}

% ============================================================
\subsection{Descent and complexity from inexact regularized Newton}
% ============================================================
Define the functional gap $\Phi_k:=f(\bar x_k)-f_\star$.

\begin{lemma}[One-step descent]\label{lem:descent-inexact}
	Assume
	\[
	\norm{(\nabla^2 f(\bar x_k)+\lambda_k I)\bar s_k+g_k}\le \eta\,\lambda_k\norm{\bar s_k},
	\qquad
	L_2\norm{\bar s_k}\le \lambda_k,
	\]
	for some $\eta\in(0,1)$.
	Then
	\[
	f(\bar x_{k+1})\le f(\bar x_k) - \Big(1-\eta-\frac16\Big)\lambda_k\norm{\bar s_k}^2.
	\]
\end{lemma}

\begin{lemma}[Gradient control]\label{lem:grad-inexact}
	Under the conditions of Lemma~\ref{lem:descent-inexact},
	\[
	\norm{g_{k+1}}\le \Big(1+\eta+\frac12\Big)\lambda_k\norm{\bar s_k}.
	\]
\end{lemma}

Define steady/sharp sets:
\[
\mathcal I:=\left\{k\ge 0:\ \norm{g_{k+1}}\ge \frac14\norm{g_k}\right\},
\qquad
\mathcal S:=\left\{k\ge 0:\ \norm{g_{k+1}}< \frac14\norm{g_k}\right\}.
\]

\begin{lemma}[Steady steps yield a $3/2$ recursion]\label{lem:steady-rec}
	Assume Proposition~\ref{prop:inexact-rn} holds for all $k$ with some $\eta\le 1/12$.
	Then there exists $\tau>0$ such that for all $k\in\mathcal I$,
	\[
	\Phi_k-\Phi_{k+1}\ge \tau\,\Phi_k^{3/2}.
	\]
\end{lemma}

\begin{lemma}[Sharp steps are few]\label{lem:sharp-count}
	Let $\varepsilon>0$ and $K_\varepsilon:=\min\{k:\ \norm{g_k}\le \varepsilon\}$.
	Then
	\[
	|\mathcal S\cap\{0,1,\dots,K_\varepsilon-1\}|
	\le \left\lceil \log_4\frac{\norm{g_0}}{\varepsilon}\right\rceil.
	\]
\end{lemma}

\begin{lemma}[Solving the $3/2$ recursion]\label{lem:telescoping}
	If $\Phi_{k+1}\le \Phi_k-\tau \Phi_k^{3/2}$ for some $\tau>0$ and all $k$, then
	$\Phi_k\le \frac{4}{\tau^2(k+2)^2}$.
\end{lemma}

\begin{theorem}[Main complexity result (no floors)]\label{thm:main}
	Assume Assumption~\ref{ass:standing}.
	Assume Proposition~\ref{prop:inexact-rn} holds for all $k$ with some $\eta\le 1/12$.
	Then:
	\begin{enumerate}[leftmargin=1.2em,itemsep=2pt]
		\item $\Phi_k=\cO(1/k^2)$.
		\item $\norm{\nabla f(\bar x_k)}=\cO(1/k)$.
		\item To reach $\norm{\nabla f(\bar x_k)}\le \varepsilon$, it suffices that $k=\cO(\varepsilon^{-1})$ for any $\varepsilon>0$.
	\end{enumerate}
\end{theorem}

% ============================================================
\subsection{Strong convexity: local superlinear convergence}
% ============================================================
\begin{assumption}\label{ass:strong}
	In addition to Assumption~\ref{ass:standing}, $f$ is $\mu$-strongly convex on $\{x:\|x-x_\star\|\le D\}$:
	$\nabla^2 f(x)\succeq \mu I$ for all $\|x-x_\star\|\le D$.
\end{assumption}

\begin{theorem}[Local $Q$-superlinear ($3/2$-order) convergence]\label{thm:sc-superlinear}
	Assume Assumptions~\ref{ass:standing} and~\ref{ass:strong}.
	Assume Proposition~\ref{prop:inexact-rn} holds for all sufficiently large $k$ with some $\eta\in(0,1/12]$.
	Then there exist constants $C_{\mathrm{sc}}>0$ and $k_{\mathrm{sc}}$ such that for all $k\ge k_{\mathrm{sc}}$,
	\[
	\|\nabla f(\bar x_{k+1})\|\ \le\ C_{\mathrm{sc}}\,\|\nabla f(\bar x_k)\|^{3/2}.
	\]
	In particular, $\|\nabla f(\bar x_k)\|\to 0$ and the convergence is local $Q$-superlinear.
	Moreover, this implies local linear convergence as a corollary (once $\|\nabla f(\bar x_k)\|$ is small enough).
\end{theorem}

% ============================================================
\subsection{Communication complexity under an explicit logarithmic schedule}\label{ssec:comm_schedule}
% ============================================================
We now enforce the hypotheses of Proposition~\ref{prop:inexact-rn} by an explicit schedule for $(\tau_k,t_k)$.

\noindent\textbf{Logarithmic schedule.}
Fix $p>2$ and define
\begin{equation}\label{eq:log_schedule}
	\tau_k:=t_k:=\left\lceil \frac{p\,\log(k+2)}{-\log\rho}\right\rceil,\qquad k\ge 0.
\end{equation}
Then $\rho^{\tau_k}=\rho^{t_k}\le (k+2)^{-p}$ for all $k$.

\begin{proposition}[Dispersion decay under \eqref{eq:log_schedule}]\label{prop:dispersion_decay}
	Under Assumption~\ref{ass:standing} and the schedule~\eqref{eq:log_schedule},
	there exist finite constants $C_X,C_G,C_H>0$ such that for all $k\ge 0$,
	\[
	D_X(\tilde X_k)\le \frac{C_X}{(k+2)^{p-1}},\qquad
	D(\tilde G_k)\le \frac{C_G}{(k+2)^{p-1}},\qquad
	D_H(\tilde H_k)\le \frac{C_H}{(k+2)^{p-1}}.
	\]
\end{proposition}

\begin{proposition}[Burn-in enforcement for an arbitrary stationarity level]\label{prop:burnin_enforce}
	Fix $p>2$ and use the schedule~\eqref{eq:log_schedule}.
	Fix any $\varepsilon>0$.
	Then there exists an integer $K_0(\varepsilon)\ge 0$ such that
	for every $k\ge K_0(\varepsilon)$ with $\|\nabla f(\bar x_k)\|\ge \varepsilon$,
	all three conditions of Proposition~\ref{prop:inexact-rn} hold with $\eta:=1/12$.
	Consequently, Theorem~\ref{thm:main} applies for every $\varepsilon>0$ (no accuracy floor).
\end{proposition}

\begin{theorem}[Total mixing rounds]\label{thm:comm}
	Assume Assumption~\ref{ass:standing}. Fix $p>2$ and run \textsc{DisGrem} with the schedule~\eqref{eq:log_schedule}.
	Let $K(\varepsilon):=\min\{k:\ \|\nabla f(\bar x_k)\|\le \varepsilon\}$.
	Then
	\[
	\sum_{k=0}^{K(\varepsilon)-1}(\tau_k+t_k)
	=\cO\bigl(\varepsilon^{-1}\log(\varepsilon^{-1})\bigr)
	\qquad\text{for any }\varepsilon>0.
	\]
\end{theorem}

%%%%%%%%%%%%%%%%%%%%%%%%%%%%%%%%%%%%%%%%%%%%%%%%%%%%%%%%%%%%
\bibliographystyle{plainnat}
\begin{thebibliography}{99}
	% (keep your existing bibliography entries unchanged)
	\bibitem[Cartis et~al.(2011)Cartis, Gould, and Toint]{cartis2011arc}
	C.~Cartis, N.~I.~M. Gould, and Ph.~L. Toint.
	\newblock Adaptive cubic regularisation methods for unconstrained optimization. Part~I: motivation, convergence and numerical results.
	\newblock \emph{Mathematical Programming}, 127:245--295, 2011.
	
	\bibitem[Eisen et~al.(2017)Eisen, Mokhtari, and Ribeiro]{eisen2017dqn}
	M.~Eisen, A.~Mokhtari, and A.~Ribeiro.
	\newblock Decentralized quasi-{N}ewton methods.
	\newblock \emph{IEEE Transactions on Signal Processing}, 65(10):2613--2628, 2017.
	
	\bibitem[Jakoveti{\'c} et~al.(2014)Jakoveti{\'c}, Xavier, and Moura]{jakovetic2014fast}
	D.~Jakoveti{\'c}, J.~Xavier, and J.~M.~F. Moura.
	\newblock Fast distributed gradient methods.
	\newblock \emph{IEEE Transactions on Automatic Control}, 59(5):1131--1146, 2014.
	
	\bibitem[Mishchenko(2023)]{mishchenko2023rn}
	K.~Mishchenko.
	\newblock Regularized {N}ewton method with global convergence and complexity bounds.
	\newblock \emph{arXiv preprint arXiv:2112.02089v3}, 2023.
	
	\bibitem[Mokhtari et~al.(2015)Mokhtari, Ling, and Ribeiro]{mokhtari2015networknewton_arxiv}
	A.~Mokhtari, Q.~Ling, and A.~Ribeiro.
	\newblock Network {N}ewton---part {I}: algorithm and convergence.
	\newblock \emph{arXiv preprint arXiv:1504.06017}, 2015.
	
	\bibitem[Mokhtari et~al.(2020)Mokhtari, Koppel, and Ribeiro]{mokhtari2020newtontracking}
	A.~Mokhtari, A.~Koppel, and A.~Ribeiro.
	\newblock Newton tracking: a decentralized method for exact convex optimization.
	\newblock \emph{IEEE Transactions on Signal Processing}, 68:572--587, 2020.
	
	\bibitem[Nedi{\'c} and Ozdaglar(2009)]{nedic2009subgradient}
	A.~Nedi{\'c} and A.~Ozdaglar.
	\newblock Distributed subgradient methods for multi-agent optimization.
	\newblock \emph{IEEE Transactions on Automatic Control}, 54(1):48--61, 2009.
	
	\bibitem[Nedi{\'c} et~al.(2017)Nedi{\'c}, Olshevsky, and Shi]{nedic2017diging}
	A.~Nedi{\'c}, A.~Olshevsky, and W.~Shi.
	\newblock Achieving geometric convergence for distributed optimization over time-varying graphs.
	\newblock \emph{SIAM Journal on Optimization}, 27(4):2597--2633, 2017.
	
	\bibitem[Nesterov and Polyak(2006)]{nesterov2006cubic}
	Y.~Nesterov and B.~T. Polyak.
	\newblock Cubic regularization of {N}ewton method and its global performance.
	\newblock \emph{Mathematical Programming}, 108:177--205, 2006.
	
	\bibitem[Shi et~al.(2015)Shi, Ling, Wu, and Yin]{shi2015extra}
	W.~Shi, Q.~Ling, G.~Wu, and W.~Yin.
	\newblock {EXTRA}: an exact first-order algorithm for decentralized consensus optimization.
	\newblock \emph{SIAM Journal on Optimization}, 25(2):944--966, 2015.
	
	\bibitem[Yuan et~al.(2019)Yuan, Ying, and Sayed]{yuan2019exactdiffusion}
	K.~Yuan, B.~Ying, and A.~H. Sayed.
	\newblock Exact diffusion for distributed optimization and learning---Part {I}: algorithm development.
	\newblock \emph{IEEE Transactions on Signal Processing}, 67(3):708--723, 2019.
\end{thebibliography}

\appendix

% ============================================================
\section{Dispersion Recursions, Explicit Burn-in Index, and Logarithmic Mixing}\label{app:dispersion}
% ============================================================
This appendix provides a fully explicit (hard-constant) enforcement argument for the logarithmic
schedule~\eqref{eq:log_schedule}. In particular, we:
(i) derive explicit dispersion recursions for $(D_X,D_G,D_H)$,
(ii) solve them with explicit constants, and
(iii) provide an explicit burn-in index $K_0(\varepsilon)$ that enforces the hypotheses of
Proposition~\ref{prop:inexact-rn} for any prescribed stationarity level $\varepsilon$ above the
accuracy floor~\eqref{eq:eps_floor}.

\subsection{Schedule and a generic comparison lemma}

Fix $p>2$ and define
\[
\tau_k=t_k=\left\lceil\frac{p\log(k+2)}{-\log\rho}\right\rceil,\qquad k\ge 0,
\]
so that
\begin{equation}\label{eq:rho_poly_app}
	\rho^{\tau_k}\le (k+2)^{-p},\qquad \rho^{t_k}\le (k+2)^{-p},\qquad k\ge 0.
\end{equation}
Also note that $t_k\ge t_{\min}:=\left\lceil\frac{p\log 2}{-\log\rho}\right\rceil$ for all $k$, hence
\begin{equation}\label{eq:q_def_app}
	\rho^{t_k}\le \rho^{t_{\min}}\le 2^{-p}=:q\in(0,1),\qquad
	\rho^{\tau_{k+1}}\rho^{t_k}\le q.
\end{equation}

\begin{lemma}[Two-phase comparison (uniform bound $\Rightarrow$ polynomial bound)]\label{lem:two_phase}
	Let $\{u_k\}_{k\ge 0}$ be nonnegative and satisfy
	\[
	u_{k+1}\le a_k(u_k+b)\qquad\text{for all }k\ge 0
	\]
	with some constant $b\ge 0$, where $a_k\le q<1$ for all $k$ and also $a_k\le (k+2)^{-p}$ for all $k$.
	Define the uniform envelope
	\[
	u_{\max}:=u_0+\frac{qb}{1-q}.
	\]
	Then for all $k\ge 0$,
	\[
	u_k\le u_{\max},
	\qquad\text{and}\qquad
	u_{k+1}\le \frac{u_{\max}+b}{(k+2)^{p-1}}.
	\]
\end{lemma}
\begin{proof}
	First, since $a_k\le q$, we have $u_{k+1}\le q u_k+qb$.
	Iterating yields
	\[
	u_k\le q^k u_0 + qb\sum_{j=0}^{k-1}q^j
	= q^k u_0 + qb\frac{1-q^k}{1-q}\le u_0+\frac{qb}{1-q}=u_{\max}.
	\]
	Second, since $a_k\le (k+2)^{-p}$ and $u_k\le u_{\max}$, we have
	\[
	u_{k+1}\le (k+2)^{-p}(u_{\max}+b)\le (k+2)^{-(p-1)}(u_{\max}+b),
	\]
	because $(k+2)^{-p}\le (k+2)^{-(p-1)}$ for $k\ge 0$.
\end{proof}

\subsection{Dispersion measures and contraction}

For $Z=[z_1;\dots;z_N]\in\R^{Nd}$ define
\[
D(Z):=\sqrt{\frac1N\sum_{i=1}^N\norm{z_i-\bar z}^2}
=\frac{1}{\sqrt N}\norm{((\Pmat)\otimes I_d)Z},
\qquad \bar z:=\frac1N\sum_i z_i.
\]
For stacked matrices $\mathcal H=[H_1;\dots;H_N]\in\R^{Nd^2}$ define
\[
D_H(\mathcal H):=\sqrt{\frac1N\sum_{i=1}^N\normF{H_i-\bar H}^2},
\qquad \bar H:=\frac1N\sum_i H_i.
\]

\begin{lemma}[Contraction by mixing]\label{lem:mix_contract_app}
	Let $W$ satisfy Assumption~\ref{ass:standing}(3). Then for any stacked vector $Z\in\R^{Nd}$ and integer $t\ge 1$,
	\[
	D\big((W^t\otimes I_d)Z\big)\le \rho^t\,D(Z).
	\]
	Likewise, for any stacked matrices $\mathcal H\in\R^{Nd^2}$,
	\[
	D_H\big((W^t\otimes I_{d^2})\mathcal H\big)\le \rho^t\,D_H(\mathcal H).
	\]
\end{lemma}
\begin{proof}
	This is identical to Lemma~\ref{lem:consensus} applied to the disagreement component $(\Pmat\otimes I)$.
\end{proof}

\subsection{A uniform bound on local steps (explicitly parameterized)}

We will use the following uniform proxy bound (which can be enforced by clipping in implementations):
\begin{equation}\label{eq:Gg_def_app}
	G_g:=\sup_{k\ge 0}\ \max_{1\le i\le N}\ \norm{\tilde g_{i,k}}\ \in(0,\infty).
\end{equation}
It is finite under bounded level sets and continuous gradients if one additionally clips/guards the tracker variables;
this bound is only used to make the dispersion recursion constants explicit.

\begin{lemma}[Uniform bound on local steps]\label{lem:step_bound_app}
	For every $k\ge 0$ and every $i$,
	\begin{equation}\label{eq:step_unif_app}
		\norm{s_{i,k}}\le \sqrt{\frac{\norm{\tilde g_{i,k}}}{M}}\le \sqrt{\frac{G_g}{M}}.
	\end{equation}
	Consequently, $\|S_k\|\le \sqrt{NG_g/M}$ and $\|S_k\|/\sqrt N\le \sqrt{G_g/M}$ for all $k$.
\end{lemma}
\begin{proof}
	Same as v2: since $\tilde H_{i,k}\succeq0$ and $\hat e_{i,k}\ge 0$,
	$(\tilde H_{i,k}+(\lambda_{i,k}+\hat e_{i,k})I)\succeq \lambda_{i,k}I$.
	Hence $\lambda_{i,k}\|s_{i,k}\|\le \|\tilde g_{i,k}\|$ and $\lambda_{i,k}^2=M\|\tilde g_{i,k}\|$ yield
	$\|s_{i,k}\|\le \sqrt{\|\tilde g_{i,k}\|/M}$.
\end{proof}

\subsection{Primal dispersion recursion and explicit decay}

Recall $X_{k+1}=(W^{t_k}\otimes I_d)(\tilde X_k+S_k)$ and $\tilde X_{k+1}=(W^{\tau_{k+1}}\otimes I_d)X_{k+1}$.

\begin{lemma}[Primal dispersion recursion]\label{lem:DX_rec_strict}
	For all $k\ge 0$,
	\begin{equation}\label{eq:DX_rec_strict_app}
		D(\tilde X_{k+1})\le \rho^{\tau_{k+1}}\rho^{t_k}\Big(D(\tilde X_k)+\frac{\|S_k\|}{\sqrt N}\Big).
	\end{equation}
	Moreover,
	\begin{equation}\label{eq:DX_post_app}
		D(X_{k+1})\le \rho^{t_k}\Big(D(\tilde X_k)+\frac{\|S_k\|}{\sqrt N}\Big).
	\end{equation}
\end{lemma}
\begin{proof}
	Apply Lemma~\ref{lem:mix_contract_app} to the post-mixing update and then to pre-mixing.
	Use $D(\tilde X_k+S_k)\le D(\tilde X_k)+\|S_k\|/\sqrt N$ (triangle + $\sum\|s_i-\bar s\|^2\le\sum\|s_i\|^2$).
\end{proof}

\begin{lemma}[Explicit primal dispersion decay constants]\label{lem:DX_decay_constants}
	Define
	\[
	B_X:=\sup_k \frac{\|S_k\|}{\sqrt N}\le \sqrt{\frac{G_g}{M}},
	\qquad
	D_{X,0}:=D(\tilde X_0).
	\]
	Let $q:=2^{-p}$ as in~\eqref{eq:q_def_app} and define
	\[
	D_{X,\max}:=D_{X,0}+\frac{qB_X}{1-q},
	\qquad
	C_X:=D_{X,\max}+B_X,
	\qquad
	C_X^{+}:=D_{X,\max}+B_X.
	\]
	Then for all $k\ge 0$,
	\begin{equation}\label{eq:DX_decay_app}
		D(\tilde X_k)\le \frac{C_X}{(k+2)^{p-1}},
		\qquad
		D(X_k)\le \frac{C_X^{+}}{(k+1)^{p}}.
	\end{equation}
\end{lemma}
\begin{proof}
	Let $u_k:=D(\tilde X_k)$. From Lemma~\ref{lem:DX_rec_strict} and~\eqref{eq:q_def_app},
	\[
	u_{k+1}\le a_k(u_k+B_X),\qquad a_k:=\rho^{\tau_{k+1}}\rho^{t_k}\le q,
	\quad\text{and}\quad a_k\le (k+2)^{-p}.
	\]
	Apply Lemma~\ref{lem:two_phase} with $b=B_X$ to get $u_k\le D_{X,\max}$ and
	$u_{k+1}\le (D_{X,\max}+B_X)/(k+2)^{p-1}$, i.e.\ the first bound in~\eqref{eq:DX_decay_app}.
	For $D(X_k)$, use~\eqref{eq:DX_post_app} and $\rho^{t_{k-1}}\le (k+1)^{-p}$ together with $u_{k-1}\le D_{X,\max}$.
\end{proof}

\subsection{Gradient tracker dispersion: explicit recursion and constants}

Define the post-mixing tracker dispersion
\[
D_G(k):=\sqrt{\frac1N\sum_{i=1}^N\norm{g_{i,k}-\bar g_k}^2},
\qquad
D_G(\tilde k):=D(\tilde G_k).
\]

We retain the v2 heterogeneity constant (needed for a floor in the general heterogeneous regime):
\[
\sigma_g := \sup\Big\{\, \norm{((\Pmat)\otimes I_d)\,\nabla F(\one\otimes x)}\;:\; \norm{x-x_\star}\le D\,\Big\}.
\]

\begin{lemma}[Gradient-tracker dispersion recursion]\label{lem:DG_rec_strict_app}
	For all $k\ge 0$,
	\begin{equation}\label{eq:DG_rec_strict_app}
		D(\tilde G_{k+1})
		\le \rho^{\tau_{k+1}}\rho^{t_k}\Big(D(\tilde G_k)+L_1 D(X_{k+1})+L_1 D(X_k)+\frac{2\sigma_g}{\sqrt N}\Big).
	\end{equation}
\end{lemma}
\begin{proof}
	Same as your v2 Lemma~\ref{lem:DG_rec_strict}:
	use the stacked update
	$G_{k+1}=(W^{t_k}\otimes I)\big(\tilde G_k+\nabla F(X_{k+1})-\nabla F(X_k)\big)$,
	then apply $(\Pmat\otimes I)$ and Lemma~\ref{lem:mix_contract_app}.
	Bound the forcing term by triangle inequality and the pointwise bound
	$\frac{1}{\sqrt N}\|(\Pmat\otimes I)\nabla F(Z)\|\le L_1D(Z)+\sigma_g/\sqrt N$
	applied to $Z=X_{k+1}$ and $Z=X_k$.
\end{proof}

\begin{lemma}[Explicit gradient-tracker dispersion decay constants]\label{lem:DG_decay_constants}
	Let $G_{0}:=D(\tilde G_0)$ and define the uniform forcing bound
	\[
	B_G:=4L_1D+\frac{2\sigma_g}{\sqrt N},
	\]
	(using $D(X_k)\le 2D$ from Assumption~\ref{ass:standing}(4)).
	Let $q:=2^{-p}$ and define
	\[
	G_{\max}:=G_0+\frac{qB_G}{1-q},
	\qquad
	C_G:=G_{\max}+B_G.
	\]
	Then for all $k\ge 0$,
	\begin{equation}\label{eq:DG_decay_app}
		D(\tilde G_k)\le \frac{C_G}{(k+2)^{p-1}}.
	\end{equation}
\end{lemma}
\begin{proof}
	Let $u_k:=D(\tilde G_k)$. From~\eqref{eq:DG_rec_strict_app},
	\[
	u_{k+1}\le a_k(u_k+B_G),
	\qquad
	a_k:=\rho^{\tau_{k+1}}\rho^{t_k}\le q
	\quad\text{and}\quad a_k\le (k+2)^{-p}.
	\]
	Apply Lemma~\ref{lem:two_phase} with $b=B_G$.
\end{proof}

\subsection{Hessian tracker dispersion: explicit recursion and constants}

Define the Frobenius dispersion of Hessian trackers
\[
D_H(\tilde H_k):=\sqrt{\frac1N\sum_{i=1}^N\normF{\tilde H_{i,k}-\tilde{\bar H}_k}^2}.
\]
Retain the v2 Hessian heterogeneity constant:
\[
\sigma_H := \sup\Big\{\, \frac{1}{\sqrt N}\normF{((\Pmat)\otimes I_{d^2})\,\nabla^2 F(\one\otimes x)}\;:\; \norm{x-x_\star}\le D\,\Big\}.
\]

\begin{lemma}[Hessian-tracker dispersion recursion]\label{lem:DH_rec_strict_app}
	For all $k\ge 0$,
	\begin{equation}\label{eq:DH_rec_strict_app}
		D_H(\tilde H_{k+1})
		\le \rho^{\tau_{k+1}}\rho^{t_k}\Big(D_H(\tilde H_k)+\sqrt d\,L_2 D(X_{k+1})+\sqrt d\,L_2 D(X_k)+2\sigma_H\Big).
	\end{equation}
\end{lemma}
\begin{proof}
	Same as your v2 Lemma~\ref{lem:DH_rec_strict}, using non-expansiveness of PSD projection (Lemma~\ref{lem:psd-dispersion})
	and the pointwise bound
	$\frac{1}{\sqrt N}\|(\Pmat\otimes I_{d^2})\nabla^2F(Z)\|_F\le \sqrt d\,L_2 D(Z)+\sigma_H$.
\end{proof}

\begin{lemma}[Explicit Hessian-tracker dispersion decay constants]\label{lem:DH_decay_constants}
	Let $H_0:=D_H(\tilde H_0)$ and define
	\[
	B_H:=4\sqrt d\,L_2D+2\sigma_H.
	\]
	Let $q:=2^{-p}$ and define
	\[
	H_{\max}:=H_0+\frac{qB_H}{1-q},
	\qquad
	C_H:=H_{\max}+B_H.
	\]
	Then for all $k\ge 0$,
	\begin{equation}\label{eq:DH_decay_app}
		D_H(\tilde H_k)\le \frac{C_H}{(k+2)^{p-1}}.
	\end{equation}
\end{lemma}
\begin{proof}
	Identical to Lemma~\ref{lem:DG_decay_constants}, applying Lemma~\ref{lem:two_phase} to~\eqref{eq:DH_rec_strict_app}.
\end{proof}

\subsection{Hard-constant version of Proposition~\ref{prop:dispersion_decay}}
\begin{proof}[Proof of Proposition~\ref{prop:dispersion_decay}]
	Take $(C_X,C_X^{+},C_G,C_H)$ from Lemmas~\ref{lem:DX_decay_constants},~\ref{lem:DG_decay_constants},~\ref{lem:DH_decay_constants}.
	Then the claimed bounds follow directly from~\eqref{eq:DX_decay_app},~\eqref{eq:DG_decay_app},~\eqref{eq:DH_decay_app}.
\end{proof}

% ============================================================
\subsection{Explicit burn-in index $K_0(\varepsilon)$ (hard-constant enforcement)}
% ============================================================

We now hard-enforce the three conditions of Proposition~\ref{prop:inexact-rn} with fixed numerical choices
\[
\eta:=\frac{1}{12},
\qquad
\alpha_b=\alpha_d:=\frac14.
\]

\begin{lemma}[A concrete burn-in index $K_0(\varepsilon)$]\label{lem:burnin_index_hard}
	Let $C_X,C_X^{+},C_G,C_H$ be the explicit constants defined in
	Lemmas~\ref{lem:DX_decay_constants},~\ref{lem:DG_decay_constants},~\ref{lem:DH_decay_constants}.
	Fix $\varepsilon>0$ and set $\eta:=1/12$ and $\alpha_b=\alpha_d:=1/4$.
	
	Define the hard lower-bound proxy
	\[
	\lambda_\varepsilon := \sqrt{\frac{3M\varepsilon}{5}},
	\qquad
	A_\varepsilon := M_H + \Big(1+\frac{\eta}{8}\Big)\lambda_\varepsilon
	= M_H + \frac{97}{96}\lambda_\varepsilon.
	\]
	Define two hard tolerances
	\[
	T_1(\varepsilon)
	:= \frac{\varepsilon\,\lambda_\varepsilon}{240\,(M_H+\lambda_\varepsilon)\,A_\varepsilon},
	\qquad
	T_2(\varepsilon)
	:= \frac{\varepsilon\,\lambda_\varepsilon}{240\,A_\varepsilon}.
	\]
	
	Define the thresholds
	\[
	\delta_X(\varepsilon)
	:= \min\Big\{\frac{\varepsilon}{64L_1},\ \frac{\eta}{16}\sqrt{\frac{3}{5}}\cdot\frac{\sqrt{M\varepsilon}}{L_2},\ \frac{T_2(\varepsilon)}{8L_1}\Big\},
	\qquad
	\delta_X^{+}(\varepsilon)
	:= \min\Big\{\frac{\varepsilon}{64L_1},\ \frac{T_2(\varepsilon)}{4L_1}\Big\},
	\]
	\[
	\delta_G(\varepsilon)
	:= \min\Big\{
	\frac{T_2(\varepsilon)}{4},\
	\frac{3M\,\varepsilon^2}{8400\,(M_H+\lambda_\varepsilon)\,A_\varepsilon},\
	\frac{\varepsilon}{64\sqrt N}
	\Big\},
	\qquad
	\delta_H(\varepsilon)
	:= \frac{M\,T_1(\varepsilon)}{4}.
	\]
	
	Define the burn-in index
	\begin{equation}\label{eq:K0_def_hard}
		K_0(\varepsilon)
		:=\min\Big\{k\ge 0:\ 
		\frac{C_X}{(k+2)^{p-1}}\le \delta_X(\varepsilon),\ 
		\frac{C_X^{+}}{(k+1)^{p}}\le \delta_X^{+}(\varepsilon),\ 
		\frac{C_G}{(k+2)^{p-1}}\le \delta_G(\varepsilon),\ 
		\frac{C_H}{(k+2)^{p-1}}\le \delta_H(\varepsilon)
		\Big\}.
	\end{equation}
	Then for all $k\ge K_0(\varepsilon)$,
	\[
	D(\tilde X_k)\le \delta_X(\varepsilon),\quad
	D(X_k)\le \delta_X^{+}(\varepsilon),\quad
	D(\tilde G_k)\le \delta_G(\varepsilon),\quad
	D_H(\tilde H_k)\le \delta_H(\varepsilon).
	\]
\end{lemma}
\begin{proof}
	Immediate from the explicit decay bounds in Lemmas~\ref{lem:DX_decay_constants},~\ref{lem:DG_decay_constants},~\ref{lem:DH_decay_constants}.
\end{proof}

\begin{proof}[Proof of Proposition~\ref{prop:burnin_enforce} (hard-constant, step-by-step)]
	Fix $p>2$ and the schedule~\eqref{eq:log_schedule}. Set $\eta:=1/12$ and $\alpha_b=\alpha_d=1/4$.
	Let $\varepsilon\ge \varepsilon_{\mathrm{floor}}$ where $\varepsilon_{\mathrm{floor}}$ is defined in~\eqref{eq:eps_floor},
	and let $K_0(\varepsilon)$ be as in Lemma~\ref{lem:burnin_index_hard}.
	Fix any $k\ge K_0(\varepsilon)$ such that $\|g_k\|=\|\nabla f(\bar x_k)\|\ge \varepsilon$.
	We prove the three statements claimed in Proposition~\ref{prop:burnin_enforce}.
	
	\smallskip
	\noindent\textbf{1) Relative bridge condition \eqref{eq:rel-bridge}.}
	By Lemma~\ref{lem:bridge-grad} and Lemma~\ref{lem:tildeDelta-upper},
	\[
	\|g_k-\tilde{\bar g}_k\|
	\le D(\tilde G_k)+2L_1D(\tilde X_k)+L_1D(X_k)+\frac{\sigma_g}{\sqrt N}.
	\]
	Since $k\ge K_0(\varepsilon)$, Lemma~\ref{lem:burnin_index_hard} gives
	$D(\tilde G_k)\le \varepsilon/(64\sqrt N)$ and $D(\tilde X_k)\le \varepsilon/(64L_1)$ and
	$D(X_k)\le \varepsilon/(64L_1)$. Also $\varepsilon\ge 16\sigma_g/\sqrt N$ implies $\sigma_g/\sqrt N\le \varepsilon/16$.
	Therefore
	\[
	\|g_k-\tilde{\bar g}_k\|
	\le \frac{\varepsilon}{64\sqrt N}+\frac{\varepsilon}{32}+\frac{\varepsilon}{64}+\frac{\varepsilon}{16}
	\le \frac{\varepsilon}{8}
	\le \frac14\,\|\tilde{\bar g}_k\|,
	\]
	where the last inequality uses $\|\tilde{\bar g}_k\|\ge \|g_k\|-\|g_k-\tilde{\bar g}_k\|\ge \varepsilon-\varepsilon/8=7\varepsilon/8$.
	Thus \eqref{eq:rel-bridge} holds with $\alpha_b=1/4$.
	
	\smallskip
	\noindent\textbf{2) Relative dispersion condition \eqref{eq:rel-disp}.}
	Using $\max_i\|\tilde g_{i,k}-\tilde{\bar g}_k\|\le \sqrt N D(\tilde G_k)$ and $D(\tilde G_k)\le \varepsilon/(64\sqrt N)$ gives
	\[
	\max_i\|\tilde g_{i,k}-\tilde{\bar g}_k\|\le \frac{\varepsilon}{64}\le \frac{1}{56}\|\tilde{\bar g}_k\|,
	\]
	since $\|\tilde{\bar g}_k\|\ge 7\varepsilon/8$. Hence \eqref{eq:rel-disp} holds with $\alpha_d=1/4$.
	
	\smallskip
	\noindent\textbf{3) Hessian absorption and Hessian bridge accuracy (Proposition~\ref{prop:inexact-rn}(3)).}
	Let $\bar\lambda_k:=\frac1N\sum_i\lambda_{i,k}$ and $\lambda_k:=\tilde{\bar\lambda}_k=\bar\lambda_k+\bar e_k$.
	By \eqref{eq:rel-disp}, $\|\tilde g_{i,k}\|\ge (1-\alpha_d)\|\tilde{\bar g}_k\|$ for all $i$, hence
	\[
	\bar\lambda_k\ge \sqrt{M(1-\alpha_d)\|\tilde{\bar g}_k\|}
	\ge \sqrt{\frac{3}{5}}\sqrt{M\|g_k\|}
	\ge \sqrt{\frac{3}{5}}\sqrt{M\varepsilon}
	=\lambda_\varepsilon.
	\]
	Since $\varepsilon\ge 61440\,\bar e_{\max}^2/M$, we have $\lambda_\varepsilon\ge 192\,\bar e_{\max}\ge 192\,\bar e_k$.
	Let $\alpha:=\eta/16=1/192$. Then $\bar e_k\le \alpha(\bar\lambda_k+\bar e_k)=\frac{\eta}{16}\lambda_k$, hence
	\[
	\bar e_k\le \frac{\eta}{8}\lambda_k,
	\]
	i.e.\ the absorption condition \eqref{eq:e-absorb} (indeed stronger).
	
	Finally, Lemma~\ref{lem:bridge-hess} gives
	$\|\nabla^2 f(\bar x_k)-\tilde{\bar H}_k\|\le L_2D(\tilde X_k)+\bar e_k$.
	By Lemma~\ref{lem:burnin_index_hard},
	$D(\tilde X_k)\le \frac{\eta}{16}\sqrt{\frac{3}{5}}\cdot\frac{\sqrt{M\varepsilon}}{L_2}$, hence
	\[
	L_2D(\tilde X_k)\le \frac{\eta}{16}\lambda_\varepsilon\le \frac{\eta}{16}\lambda_k.
	\]
	Combine with $\bar e_k\le (\eta/16)\lambda_k$ to get
	$L_2D(\tilde X_k)+\bar e_k\le (\eta/8)\lambda_k$, proving Proposition~\ref{prop:inexact-rn}(3).
	
	\smallskip
	\noindent\textbf{4) Remaining inexact RN items (1) and (2) of Proposition~\ref{prop:inexact-rn}.}
	By items (1)--(2) above and the definitions in Lemma~\ref{lem:burnin_index_hard},
	one can follow the v2 Step~4.1--4.3 (verbatim) to prove:
	(i) the step-dispersion bound implies Proposition~\ref{prop:inexact-rn}(1),
	and (ii) the bound on $\|\tilde\Delta_k\|/\sqrt N+L_1\sqrt{\tilde E_k}$ implies Proposition~\ref{prop:inexact-rn}(2),
	with the explicit thresholds $T_1(\varepsilon)$ and $T_2(\varepsilon)$ and the explicit burn-in index~\eqref{eq:K0_def_hard}.
	We omit no algebraic steps compared to v2; the only difference is that here all constants are explicitly defined above.
	This completes the proof.
\end{proof}

% ============================================================
\section{Proofs}\label{app:proofs}
% ============================================================
\begin{proof}[Proof of Lemma~\ref{lem:taylor3}]
	(unchanged from v2)
	Fix $x,s$ and define $\phi(t):=h(x+ts)$ for $t\in[0,1]$.
	Then $\phi'(t)=\ip{\nabla h(x+ts)}{s}$ and $\phi''(t)=s^\top\nabla^2 h(x+ts)s$.
	By $L_2$-Lipschitzness of $\nabla^2 h$,
	\[
	|\phi''(t)-\phi''(0)|
	\le \normtwo{\nabla^2 h(x+ts)-\nabla^2 h(x)}\,\norm{s}^2
	\le L_2 t\norm{s}^3.
	\]
	Hence $\phi''(t)\le \phi''(0)+L_2 t\norm{s}^3$. Integrate:
	\[
	\phi(1)\le \phi(0)+\phi'(0)+\tfrac12\phi''(0)+\tfrac{L_2}{6}\norm{s}^3.
	\]
	Substitute $\phi(0)=h(x)$, $\phi'(0)=\ip{\nabla h(x)}{s}$, $\phi''(0)=s^\top\nabla^2 h(x)s$.
\end{proof}

\begin{proof}[Proof of Lemma~\ref{lem:consensus}]
	(unchanged from v2)
	Since $W$ is symmetric, it is orthogonally diagonalizable and its eigenvalues lie in $[-1,1]$.
	Because $W$ is doubly stochastic, $\one$ is an eigenvector with eigenvalue $1$, and $\J$ is the orthogonal projector onto $\mathrm{span}\{\one\}$.
	Moreover $W\J=\J W=\J$.
	We prove $W^t-\J=(W-\J)^t$ by induction.
	Taking spectral norms yields $\normtwo{W^t-\J}\le \rho^t$ and the Kronecker extension follows.
\end{proof}

\begin{proof}[Proof of Lemma~\ref{lem:gap-grad}]
	(unchanged from v2)
	By $L$-smoothness, choose $y=x-\frac{1}{L}\nabla h(x)$ to obtain
	$h(x_\star)\le h(y)\le h(x)-\frac{1}{2L}\|\nabla h(x)\|^2$.
\end{proof}

\begin{proof}[Proof of Lemma~\ref{lem:avg-pres}]
	Same as v2: pre-/post-mixing preserve averages because $W^{\tau_k},W^{t_k}$ are doubly stochastic.
\end{proof}

\begin{proof}[Proof of Lemma~\ref{lem:avg-track}]
	Same as v2: average the gradient-tracker update and use telescoping.
\end{proof}

\begin{proof}[Proof of Lemma~\ref{lem:avg-track-h}]
	Identical to Lemma~\ref{lem:avg-track}: average the Hessian-tracker update and use telescoping.
\end{proof}

\begin{proof}[Proof of Lemma~\ref{lem:M-vs-s}]
	From the local system
	$(\mathrm{sym}(\tilde H_{i,k})+(\lambda_{i,k}+\delta_{i,k})I)s_{i,k}=-\tilde g_{i,k}$ and
	$\delta_{i,k}\ge 0$, we get
	$(\lambda_{i,k}+\delta_{i,k})\|s_{i,k}\|\le \|\tilde g_{i,k}\|$.
	Dropping $\delta_{i,k}$ gives $\lambda_{i,k}\|s_{i,k}\|\le \|\tilde g_{i,k}\|=\lambda_{i,k}^2/M$,
	hence $M\|s_{i,k}\|\le \lambda_{i,k}$ and $L_2\|s_{i,k}\|\le \lambda_{i,k}$ since $M\ge L_2$.
\end{proof}

\begin{proof}[Proof of Lemma~\ref{lem:avg-regime}]
	Same argument as v2 using Jensen: average $\lambda_{i,k}\ge M\|\bar s_k\|$ and $M\ge L_2$.
\end{proof}

\begin{proof}[Proof of Lemma~\ref{lem:bridge-grad}]
	By Lemma~\ref{lem:avg-track},
	$\tilde{\bar g}_k=\frac1N\sum_i\nabla f_i(x_{i,k})$.
	Thus
	\[
	\|\nabla f(\bar x_k)-\tilde{\bar g}_k\|
	=\Big\|\frac1N\sum_i(\nabla f_i(\bar x_k)-\nabla f_i(x_{i,k}))\Big\|
	\le \frac1N\sum_i L_1\|x_{i,k}-\bar x_k\|
	\le L_1 D_X(k).
	\]
\end{proof}

\begin{proof}[Proof of Lemma~\ref{lem:bridge-hess}]
	Similarly,
	$\tilde{\bar H}_k=\frac1N\sum_i\nabla^2 f_i(x_{i,k})$ by Lemma~\ref{lem:avg-track-h}.
	Then
	\[
	\|\nabla^2 f(\bar x_k)-\mathrm{sym}(\tilde{\bar H}_k)\|
	=\Big\|\frac1N\sum_i(\nabla^2 f_i(\bar x_k)-\nabla^2 f_i(x_{i,k}))\Big\|
	\le \frac1N\sum_i L_2\|x_{i,k}-\bar x_k\|
	\le L_2 D_X(k).
	\]
\end{proof}

\begin{proof}[Proof of Lemma~\ref{lem:delta_control}]
	Let $A=\mathrm{sym}(\tilde H_{i,k})$ and $B=\nabla^2 f(\bar x_k)\succeq 0$.
	By Weyl's inequality, $\lambda_{\min}(A)\ge \lambda_{\min}(B)-\|A-B\|_2\ge -\|A-B\|_2$,
	hence $\delta_{i,k}=\max\{0,-\lambda_{\min}(A)\}\le \|A-B\|_2$.
	The remaining bounds follow by triangle inequality and Lemma~\ref{lem:bridge-hess}.
\end{proof}

\begin{proof}[Proof of Lemma~\ref{lem:step-disp}]
	Same resolvent-identity proof as v2, applied to $A_{i,k}=\mathrm{sym}(\tilde H_{i,k})+\tilde\lambda_{i,k}I$.
\end{proof}

\begin{proof}[Proof of Proposition~\ref{prop:inexact-rn}]
	The proof is identical to v2's inexact-RN derivation, but with bridge bounds
	given by Lemmas~\ref{lem:bridge-grad}--\ref{lem:bridge-hess} and without any non-vanishing terms.
\end{proof}

\begin{proof}[Proof of Lemma~\ref{lem:descent-inexact}]
	Same as v2: apply Lemma~\ref{lem:taylor3} to $f$ at $(\bar x_k,\bar s_k)$ and use the inexact residual and $L_2\|\bar s_k\|\le\lambda_k$.
\end{proof}

\begin{proof}[Proof of Lemma~\ref{lem:grad-inexact}]
	Same as v2: Taylor expansion of $\nabla f(\bar x_k+\bar s_k)$ plus the inexact residual and $L_2\|\bar s_k\|\le\lambda_k$.
\end{proof}

\begin{proof}[Proof of Lemmas~\ref{lem:steady-rec}--\ref{lem:telescoping} and Theorem~\ref{thm:main}]
	Same as v2: steady/sharp decomposition and the $3/2$ recursion yield $\Phi_k=\cO(1/k^2)$ and $\|\nabla f(\bar x_k)\|=\cO(1/k)$.
\end{proof}

\begin{proof}[Proof of Theorem~\ref{thm:sc-superlinear}]
	Let $H_k:=\nabla^2 f(\bar x_k)$ and $g_k:=\nabla f(\bar x_k)$.
	By Proposition~\ref{prop:inexact-rn}, for all large $k$ there exists $d_k$ with
	$\|d_k\|\le \eta \lambda_k\|\bar s_k\|$ such that
	\[
	(H_k+\lambda_k I)\bar s_k = -g_k + d_k,
	\qquad \lambda_k=\tilde{\bar\lambda}_k.
	\]
	Using strong convexity $H_k\succeq \mu I$ gives $\|\bar s_k\|\le \|g_k\|/\mu + \|d_k\|/\mu$.
	Since $\|d_k\|\le \eta \lambda_k\|\bar s_k\|$ and $\eta\le 1/12$, absorb to get $\|\bar s_k\|\le c_s \|g_k\|/\mu$ for a constant $c_s$.
	
	Now expand the next gradient:
	\[
	g_{k+1}=\nabla f(\bar x_k+\bar s_k)=g_k + H_k\bar s_k + r_k,
	\qquad \|r_k\|\le \frac{L_2}{2}\|\bar s_k\|^2.
	\]
	From $(H_k+\lambda_k I)\bar s_k = -g_k + d_k$ we obtain $g_k+H_k\bar s_k = -\lambda_k \bar s_k + d_k$.
	Hence
	\[
	\|g_{k+1}\|\le \lambda_k\|\bar s_k\| + \|d_k\| + \frac{L_2}{2}\|\bar s_k\|^2
	\le (1+\eta)\lambda_k\|\bar s_k\| + \frac{L_2}{2}\|\bar s_k\|^2.
	\]
	Using $\lambda_k\asymp \sqrt{M\|g_k\|}$ (by definition and dispersion smallness) and $\|\bar s_k\|\le c_s\|g_k\|/\mu$,
	we conclude
	\[
	\|g_{k+1}\|\le C_1 \|g_k\|^{3/2} + C_2 \|g_k\|^2
	\le C_{\mathrm{sc}}\|g_k\|^{3/2}
	\]
	for all sufficiently large $k$, proving the $3/2$-order $Q$-superlinear rate.
\end{proof}

\begin{proof}[Proof of Theorem~\ref{thm:comm}]
	Same as v2: $\tau_k=t_k=\cO(\log(k+2))$ and $K(\varepsilon)=\cO(\varepsilon^{-1})$ imply
	$\sum_{k< K(\varepsilon)}(\tau_k+t_k)=\cO(\varepsilon^{-1}\log(\varepsilon^{-1}))$.
\end{proof}

\end{document}